\section{Theory}
\label{sec:theory}

\subsection{Assumptions}
\label{sec:assumptions}

\begin{assumption}[Sampling and error]
\label{ass:sample}
The observations $(Y_i,\bX_i,\bU_i)$, $i \le n$, are independent and
identically distributed and follow \eqref{eq:model}, and the error is
conditionally sub-Gaussian: there is $\sigma < \infty$ with
$\E\{\exp(t\varepsilon_i) \mid \bX_i,\bU_i\} \le \exp(\sigma^2t^2/2)$ almost
surely, for all real $t$.
\end{assumption}

\begin{assumption}[Linear covariates]
\label{ass:X}
$\max_{\cv}|X_{\cv}| \le B_X$ almost surely, and there are constants
$0 < \kappa_1 \le \kappa_2 < \infty$ with
$\kappa_1 \le \lmin\{\E(\bX\bX\T\mid\bU)\} \le
\lmaxx\{\E(\bX\bX\T\mid\bU)\} \le \kappa_2$ almost surely.
\end{assumption}

\begin{assumption}[Modulating covariates]
\label{ass:U}
$\bU$ takes values in $[0,1]^q$ and has a joint Lebesgue density $p_{\bU}$
with $0 < c_U \le p_{\bU}(u) \le C_U < \infty$ for every $u \in [0,1]^q$.
\end{assumption}

\begin{assumption}[Basis]
\label{ass:basis}
The wavelets are the periodization \eqref{eq:periodized} of an orthonormal,
compactly supported family with $\psi$ bounded and $N$ vanishing moments, and
$\jzero = 0$.
\end{assumption}

\begin{assumption}[Besov regularity]
\label{ass:besov}
There are $s > 0$, $1 \le \pi, r \le \infty$ and $C_g < \infty$ such that
every component satisfies $\int_0^1 \gcomp{\cv}{\md} = 0$ and, writing
$\coefs{\cv}{\md}{\lev}{\tsl} = \inner{\gcomp{\cv}{\md}}{\wav{\lev}{\tsl}}$
and $\theta^{*}_{\cv\md,\lev\cdot}$ for the vector of the coefficients of
level $\lev$,
\[
  \norm{\theta^{*}_{\cv\md,\lev\cdot}}_{\pi} \;\le\;
  C_g\,2^{-\lev(s + 1/2 - 1/\pi)} \qquad \text{for every } \lev \ge 0 .
\]
The effective regularity $\seff = s - (1/\pi - 1/2)_{+}$ is positive.
\end{assumption}

Assumptions~\ref{ass:X} and \ref{ass:U} are the quantitative forms of the two
conditions of Proposition~\ref{prop:identif}. In the second it is the joint
density that is bounded below, not the marginals: if $U_2 = U_1$ almost surely
both marginals are uniform and the design is singular.
Assumption~\ref{ass:besov} is taken as a condition on the sequence of
coefficients rather than as membership in a function space, as in
\citet{Klopp-Pensky-2015}; when $\psi$ has more than $s$ vanishing moments and
enough Hölder regularity the two are equivalent, on the circle
\citep[Section 3.6]{Cohen-2003}, and the implication that the proofs use needs
only the vanishing moments. The index $\seff$ is the regularity that linear
approximation sees and $s$ is the one that nonlinear approximation sees; they differ
exactly when $\pi < 2$, which is the regime of inhomogeneous smoothness.

\subsection{The main result}
\label{sec:main}

Write $\widehat{f}_i = \Amat_i\chat + \Bmat_i\thetahat$ for the fit, $\Gcal_{n}$ for the
partition of Section~\ref{sec:estimator} at resolution $J_n$ with $b = \bn$,
and fix a confidence
level $\alpha \in (0,1)$. The penalty is the one that the oracle inequality of
Section~\ref{sec:oracle} calls for, namely, with $\gamma = \kappa_1c_U$
and $\Lambda = \kappa_2C_U + \gamma/2$ the constants of
Proposition~\ref{prop:gram} below,

% derivations/08-blocos.tex, Lema 14(ii)-(iii) e a Hipotese B.

\begin{equation}
  \label{eq:lambda}
  \lamG \;=\; 2\max_{G\in\Gcal_{n}}\frac{\lamzbar{G}}{w_{G}},
  \quad
  \lamzbar{G} \;=\; \frac{\sigma}{\sqrt{n}}\Bigl\{B_X\sqrt{2C_U|G|}
  + \sqrt{2\Lambda\log(|\Gcal_{n}|/\alpha)}\Bigr\} .
\end{equation}
For every chunk, $\lamG w_{G} \le \varrho\lamGone$, where
$\lamGone = 2\max_{G}\lamzbar{G}$ is of order $\sigma(\log n/n)^{1/2}$, the
order of the penalty of the lasso.

% derivations/05-taxas.tex, Corolario 5 (numeracao global); D16. Substituido
% em k = 4 pelo Corolario 11 de derivations/08-blocos.tex (D47(a)).

% derivations/08-blocos.tex, Corolario 11 (numeracao global); D47(a).

\begin{theorem}[Rate of nonlinear approximation]
\label{thm:main}
Let Assumptions~\ref{ass:sample} to \ref{ass:besov} hold with $p$, $q$ and all
the constants fixed as $n$ grows, and let the weights have ratio
$\varrho \le \bar\varrho$ for a constant $\bar\varrho$ fixed as $n$ grows
($\bar\varrho^{2} = 3$ for the weights of Section~\ref{sec:estimator}). Suppose
$\seff > s/(2s+1)$ and take $J_n = \lceil c \log_2 n\rceil$ with
\begin{equation}
  \label{eq:window}
  \frac{s}{(2s+1)\seff} \;\le\; c \;<\; 1 ,
\end{equation}
a window that is nonempty precisely under that condition. Then, with
$\pen = \lamG$ as in \eqref{eq:lambda} and
$r_n = n^{-2s/(2s+1)}(\log n)^{(2/\pi-1)_{+}/(2s+1)}$,
\[
  \normn{\widehat{f}-f}^{2} \;=\; O_p(r_n),
  \qquad
  \sum_{\cv,\md}\bigl\lVert \widehat{g}_{\cv\md}-\gcomp{\cv}{\md}
  \bigr\rVert_{L_2[0,1]}^{2} \;=\; O_p(r_n) .
\]
In particular, when $\pi \ge 2$ the rate is $n^{-2s/(2s+1)}$, with no
logarithmic factor.
\end{theorem}


Three readings of Theorem~\ref{thm:main} matter. First, the exponent is the
one of nonlinear wavelet approximation, $2s/(2s+1)$, and not the one of
linear approximation, $2\seff/(2\seff+1)$ of Theorem~\ref{thm:sieve} below; the two
agree when $\pi \ge 2$ and separate when $\pi < 2$, which is where the
 penalty pays for itself. Second, the rate is attained under
Assumption~\ref{ass:besov} alone:  no condition is imposed
on the coefficients on top of it, because the Besov condition already bounds
the ideal risk over chunks that governs the estimation error
(Lemma~\ref{lem:weak}). 
% E5g (D57, pergunta 51(b)): a terceira leitura passa a citar a cota inferior
% do Theorem S6.1 do supp (Teorema 5 de derivations/09-cota-inferior.tex).
Third,  the rate
cannot be improved. With Gaussian errors and any law of $(\bX,\bU)$ that
satisfies Assumptions~\ref{ass:X} and~\ref{ass:U}, the minimax risk over the
class of Assumption~\ref{ass:besov}, in the prediction norm and for the
components, is at least a constant multiple of $pq\,n^{-2s/(2s+1)}$, for every
$\pi$ (Theorem~S6.1 in the Supplementary Material). This is the regression
counterpart of the lower bound of \citet[Theorem~5]{Donoho-Johnstone-1998} for
Besov bodies in the Gaussian sequence model, and with a fixed number of
covariates it extends that of \citet[Theorem~1]{Klopp-Pensky-2015}, who have a
single modulator independent of the covariates. Theorem~\ref{thm:main} is
therefore minimax optimal when $\pi \ge 2$, with any number of modulators and
with covariates that depend on them. When
$\pi < 2$ the remaining factor $(\log n)^{(2/\pi-1)/(2s+1)}$ is the one in the
rate of \citet[eq.~3.19]{Klopp-Pensky-2015} with a fixed number of covariates,
and in that of block thresholding in the sequence model
\citep[Theorem~4]{Cai-1999};  the lower bound carries
no such factor, and whether the factor is necessary for this estimator is not
known. Compared with the
coordinatewise lasso, whose rate under the same conditions is
$(\log n/n)^{2s/(2s+1)}$ (Theorem~S8.3 in the Supplementary Material), the
chunks remove a factor $(\log n)^{2\seff/(2s+1)}$.
% D55 (pergunta 49(a)): a frase de D49(c) passa a valer para todo pi, pela S6.2.
 When $s < 1/2$, the prediction
rate of Theorem~\ref{thm:main} holds with no design condition at all, under
its other assumptions and in the same window for $J_{n}$ (Proposition~S6.2
in the Supplementary Material).

The proof combines three ingredients, given below: the approximation error of
$\WJ$, the design condition, and a nonasymptotic oracle inequality.
Section~\ref{sec:rates} assembles them, first at the oracle, which gives the rate of linear approximation, and then at a
comparator that keeps only the chunks of large energy,
which gives Theorem~\ref{thm:main}.

\subsection{Approximation and the design condition}
\label{sec:approx}

% derivations/02-aproximacao-besov.tex, Lema 2 e Corolario 1 (numeracao global).
\begin{lemma}[Approximation bias]
\label{lem:bias}
Under Assumptions~\ref{ass:basis} and \ref{ass:besov}, for every pair
$(\cv,\md)$ and every $J \ge 0$,
\[
  \bigl\lVert \gcomp{\cv}{\md} - \PiJ\gcomp{\cv}{\md}
  \bigr\rVert_{L_2[0,1]}^{2}
  \;=\; \sum_{\lev \ge J}\norm{\theta^{*}_{\cv\md,\lev\cdot}}_{2}^{2}
  \;\le\; \frac{C_g^{2}}{1-2^{-2\seff}}\,2^{-2J\seff} .
\]
If in addition Assumptions~\ref{ass:X} and \ref{ass:U} hold, then, with
$f_J(x,u) = \sum_{\cv}x_{\cv}\{\lvl{\cv} +
\sum_{\md}(\PiJ\gcomp{\cv}{\md})(u_{\md})\}$,
\begin{equation}
  \label{eq:bias}
  \bigl\lVert f - f_J\bigr\rVert_{L_2(P)}^{2} \;\le\;
  \frac{(pq)^{2}B_X^{2}C_UC_g^{2}}{1-2^{-2\seff}}\,2^{-2J\seff}
  \;=:\; \mathcal{B}_{J}^{2} .
\end{equation}
\end{lemma}

The exponent is $\seff$ and not $s$, and this is where the loss
$(1/\pi-1/2)_{+}$ enters: inside a level, an $\ell_\pi$ bound with $\pi < 2$
controls the $\ell_2$ norm without gain. The bias uses only the bound on the
linear covariates and the marginal density of each modulator, so dependence
between $\bX$ and $\bU$ does not affect it. What it does affect is the passage
from the prediction norm to the norm of each component, which is the design
condition.

Let $\bPsi(u) \in \R^{K}$, $K = 1 + q\NJ$, collect the constant and the
wavelets of the $q$ modulators, so that the row of $\bZ$ in
\eqref{eq:design} is a permutation of the Kronecker product
$\bX \otimes \bPsi(\bU)$. Write $\Gram = \E(\bZ\bZ\T)$ and
$\Gramhat = \bZ\T\bZ/n$ for the population and empirical Gram matrices and
$\Gram_{\bPsi} = \E\{\bPsi(\bU)\bPsi(\bU)\T\}$ for the part in the modulators
alone.

% derivations/03-desenho-produtos.tex, Proposicao 3 (numeracao global).
\begin{proposition}[Gram matrix of the product design]
\label{prop:gram}
Let Assumptions~\ref{ass:X} to \ref{ass:basis} hold and write $D = p + d$ and
$\gamma = \kappa_1 c_U$.
\begin{enumerate}
\item[(i)] $c_U \le \lmin(\Gram_{\bPsi}) \le \lmaxx(\Gram_{\bPsi}) \le C_U$
for every $J$, and consequently
$\gamma \le \lmin(\Gram) \le \lmaxx(\Gram) \le \kappa_2 C_U$. If $\bX$ and
$\bU$ are independent, $\Gram$ is, up to the permutation of columns, the
Kronecker product $\E(\bX\bX\T)\otimes\Gram_{\bPsi}$, and the two bounds are
attained as products of eigenvalues.
\item[(ii)] With $R_J = pB_X^{2}(1 + qC_\psi 2^{J})$, where $C_\psi$ depends
only on the wavelet family, $\norm{\bZ_i}_2^{2} \le R_J$ almost surely and
\[
  \Prob\Bigl(\lmin(\Gramhat) \ge \frac{\gamma}{2}\Bigr)
  \;\ge\; 1 - 2D\exp\Bigl\{-\frac{n\gamma^{2}}
  {8R_J(\kappa_2C_U + \gamma/3)}\Bigr\},
\]
which tends to one whenever $pq2^{J}\log(pq2^{J})/n \to 0$.
\end{enumerate}
\end{proposition}

Part (i) is stated for the general case: independence between the covariates
and the modulators buys the exact factorization but not a better constant. A
lower bound on the smallest eigenvalue dominates every compatibility and
restricted eigenvalue constant, including the versions for groups
of \citet{Lounici-Pontil-vandeGeer-Tsybakov-2011}, on every support and with no cone restriction,
since $\delta\T A \delta \ge \lmin(A)\norm{\delta}_2^2 \ge
\lmin(A)\norm{\delta_S}_1^2/|S|$, and likewise
$\delta\T A\delta \ge \lmin(A)\{\sum_{G\in\Hcal}w_{G}\norm{\delta_{G}}_{2}\}^{2}/\Wcal(\Hcal)$
for a set $\Hcal$ of chunks, with $\Wcal(\Hcal) = \sum_{G\in\Hcal}w_{G}^{2}$,
and it is in this form that the next
subsection consumes it. In the resolution regime of Section~\ref{sec:rates}
the requirement of part (ii) is the
rate times $\log n$, which tends to zero, so it
costs nothing there; if $pq$ were larger than $n^{1/2}$, only the
compatibility constant could be transferred, through the entrywise deviation
\citep[Corollary 6.8]{buhlmann2011statistics}, at the price of a factor
$s_0^{2}$ and of the cone.

\subsection{A nonasymptotic oracle inequality}
\label{sec:oracle}

Let $\betastar = (\bc,\thetastar)$ be the oracle of the approximation space, that is the vector
whose components are the levels and the coefficients
$\coefs{\cv}{\md}{\lev}{\tsl}$ of Assumption~\ref{ass:besov} for $\lev < J$,
so that $f_J$ is its fit, and let


\[
  \bb = f - f_J \in \R^{n},
  \qquad
  \bbar = f - \bZ\bar\bbeta, \qquad \vbar = \thetahat - \thbar ,
  \qquad \Gbar = \{G : \thbar_{G} \ne 0\},
\]
for a comparator $\bar\bbeta = (\bar\bc,\thbar) \in \R^{p}\times\R^{d}$, any
vector of levels and coefficients, $\Gbar$ being its active chunks, and
$\bar\Wcal = \Wcal(\Gbar) = \sum_{G\in\Gbar}w_{G}^{2}$.
Let $\PA$ be the orthogonal projection onto the column space of $\Amat$, which
exists whenever $\Amat$ has rank $p$, let $\MA = I_n - \PA$, and let
$\Btil = \MA\Bmat$ and $\gtil = \lmin(\Btil\T\Btil/n)$ be the residualized
design and its smallest eigenvalue, and
$\Psitil_{G} = \Btil_{G}\T\Btil_{G}/n$ the residualized Gram matrix of the
columns of a chunk. Finally let
$\smax^{2} = \max_a (\Gramhat)_{aa}$ be the largest empirical norm of a
penalized column.

% derivations/04-oraculo.tex, Teorema 1 (numeracao global). Substituido em
% k = 4 pelo Teorema 3 de derivations/08-blocos.tex; o do lasso e o
% Theorem S8.1 do supp.

% derivations/08-blocos.tex, Teorema 3 e Lema 14(i) (numeracao global).

\begin{theorem}[Oracle inequality]
\label{thm:oracle}
Let Assumption~\ref{ass:sample} hold, suppose $\Amat$ has rank $p$, let the
weights not depend on the responses, and, for $\alpha \in (0,1)$, let
\[
  \lamz{G} \;=\; \frac{\sigma}{\sqrt{n}}\Bigl[
  \{\mathrm{tr}(\Psitil_{G})\}^{1/2}
  + \bigl\{2\opnorm{\Psitil_{G}}\log(|\Gcal|/\alpha)\bigr\}^{1/2}\Bigr]
\]
and $\pen_{w} = 2\max_{G\in\Gcal}\lamz{G}/w_{G}$, and take $\pen \ge \pen_{w}$. On the event
$\Tw = \{\max_{G}\norm{(\Btil\T\bveps/n)_{G}}_{2}/w_{G} \le \pen/2\}$, which
has probability at least $1 - \alpha$, the following hold simultaneously for
every comparator $\bar\bbeta$:
\begin{enumerate}
\item[(i)] with no condition on the design,
$\tfrac12\normn{\Btil\vbar}^{2} + \pen\nGw{\thetahat} \le
3\pen\nGw{\thbar} + 2\normn{\bbar}^{2}$;
\item[(ii)] if $\gtil > 0$,
\[
  \normn{\Btil\vbar}^{2} \le \frac{64\pen^{2}\bar\Wcal}{\gtil} + 16\normn{\bbar}^{2},
  \qquad
  \norm{\vbar}_{2}^{2} \le \frac{64\pen^{2}\bar\Wcal}{\gtil^{2}}
  + \frac{16\normn{\bbar}^{2}}{\gtil} ;
\]
\item[(iii)] in either regime,
$\normn{\widehat{f}-f} \le \normn{\Btil\vbar} + 2\normn{\bbar}
+ \normn{\PA\bveps}$, and
$\E(\normn{\PA\bveps}^{2}\mid\bX,\bU) \le \sigma^{2}p/n$.
\end{enumerate}
\end{theorem}


Four features separate Theorem~\ref{thm:oracle} from the standard treatment
\citep[Chapter 6]{buhlmann2011statistics} and from that of the group
lasso by \citet{Lounici-Pontil-vandeGeer-Tsybakov-2011}. The unpenalized levels are not
placed inside the support, which would require a compatibility condition on a
cone; they are profiled out, and the Schur complement transfers the full
smallest eigenvalue of $\Gramhat$ to the residualized block,
$\gtil \ge \lmin(\Gramhat)$, so that Proposition~\ref{prop:gram}(ii) is used
as it stands and the constant does not depend on $\Gbar$. The model is
misspecified inside the approximation space, and the deterministic residual $\bbar$ is handled
by Young's inequality rather than by duality, which keeps the statement
nonasymptotic and valid also when the bias dominates, and the
statement holds for every comparator at once, which is what
Section~\ref{sec:rates} uses.  The penalty is calibrated chunk by
chunk, by the tail inequality for quadratic forms of
\citet[Theorem~2.1]{Hsu-Kakade-Zhang-2012} with a union bound over the chunks
(Lemma~S5.2 in the Supplementary Material); the event $\Tw$ does not depend on
the weights, which only choose the chunk that dictates $\pen_{w}$, and the
price of unequal weights is the factor $\varrho$ in
$\pen_{w}w_{G} \le 2\varrho\max_{G'}\lamz{G'}$, paid in the estimation term.
Since $\mathrm{tr}(\Psitil_{G}) \le |G|\smax^{2}$ and
$\opnorm{\Psitil_{G}} \le \lmaxx(\Gramhat)$, the deterministic level
\eqref{eq:lambda} dominates $\pen_{w}$ whenever $\smax^{2} \le 2B_X^2C_U$ and
$\lmaxx(\Gramhat) \le \Lambda$, two events of probability tending to one in
the regime of Section~\ref{sec:rates}; it is the empirical column norms that
enter, and not $\max_{i,a}|Z_{ia}|$, which is of order $2^{J/2}$ and would
destroy the rates. The software standardizes the columns and orthonormalizes
each chunk before penalizing it; the theorem covers that estimator as well,
with $\Psitil_{G}$ replaced by $Q_{G}^{-1/2}\Psitil_{G}Q_{G}^{-1/2}$, $Q_{G}$
being the centered Gram matrix of the chunk, which makes $\lamz{G}$ pivotal,
and with $\bar\Wcal$ multiplied by $\max_{G}\lmaxx(Q_{G})$ (Remark~S5.2 in the
Supplementary Material). The term $\sigma^{2}p/n$ in part (iii) is the price
of leaving the $p$ levels out of the penalty; no choice of $\pen$ removes it,
it is of smaller order than the estimation term whenever
$p \lesssim |\Gbar|\log n$, and it is absent from
\citet{Klopp-Pensky-2015}, where the levels are penalized along with
everything else.

\subsection{Rates}
\label{sec:rates}

Under Assumption~\ref{ass:besov} the oracle is generically dense: the
coefficients decay but do not vanish, so that every chunk is
active and part (ii) of
Theorem~\ref{thm:oracle}, read at the oracle, delivers the rate of
linear approximation, now with no logarithmic factor.

% derivations/05-taxas.tex, Teorema 2 (numeracao global). Substituido em k = 4
% pelo Teorema 4 de derivations/08-blocos.tex; o do lasso e o Theorem S8.2.

% derivations/08-blocos.tex, Teorema 4 (numeracao global).

\begin{theorem}[Rate of linear approximation]
\label{thm:sieve}
Let Assumptions~\ref{ass:sample} to \ref{ass:besov} hold with $p$, $q$ and the
constants fixed, let the weights be as in Theorem~\ref{thm:main}, and take
\begin{equation}
  \label{eq:Jn}
  J_n = \min\bigl\{J \in \mathbb{N} :\;
  2^{J} \ge n^{1/(2\seff+1)}\bigr\}.
\end{equation}
Then the requirement of Proposition~\ref{prop:gram}(ii) holds along $J_n$,
being of order $n^{-2\seff/(2\seff+1)}\log n$, and, with $\pen = \lamG$ as in
\eqref{eq:lambda},
\[
  \normn{\widehat{f}-f}^{2}
  \;=\; O_p\bigl(n^{-2\seff/(2\seff+1)}\bigr) .
\]
At this resolution the estimation term, of order $pq(\bn + 2^{J_n})/n$, and
the bias term are of the same order, and the term $\sigma^{2}p/n$ of
Theorem~\ref{thm:oracle}(iii) is of strictly smaller order.
\end{theorem}


% derivations/05-taxas.tex, Corolario 4 (numeracao global). Substituido em
% k = 4 pelo Corolario 10 de derivations/08-blocos.tex.

% derivations/08-blocos.tex, Corolario 10 (numeracao global).

\begin{corollary}[Components and coefficients]
\label{cor:components}
In the setting of Theorem~\ref{thm:sieve}, with
$\rhoG = n^{-\seff/(2\seff+1)}$,
\[
  \sum_{\cv,\md}\bigl\lVert\widehat{g}_{\cv\md}-\gcomp{\cv}{\md}
  \bigr\rVert_{L_2[0,1]}^{2} = O_p\bigl\{(\rhoG)^{2}\bigr\},
  \qquad
  \bigl\lVert\widehat{\fcoef{\cv}}-\fcoef{\cv}
  \bigr\rVert_{L_2([0,1]^{q})} = O_p(\rhoG)
\]
for $\cv = 1,\dots,p$, and $\norm{\chat - \bc}_{2} = O_p(\rhoG)$.
\end{corollary}


The passage from the prediction norm to the norm of a single component is
where the design condition is used a second time, and the smallest eigenvalue
appears squared in the constant. In the arguments of this
paper the design condition is needed for the components in every case, and
for the prediction rate of Theorem~\ref{thm:main} only when $s \ge 1/2$.
% E5f (E1.15; D54): a frase acima qualifica a condicao de desenho; a
% introducao (D50) nao muda.
Two further statements complete the picture.
The first dispenses with the design condition altogether, at the price of a
slower rate when $s \ge 1/2$, and is the unconditional statement of this paper.

% derivations/05-taxas.tex, Proposicao 4 (numeracao global). Substituida em
% k = 4 pelo Corolario 14 de derivations/08-blocos.tex (D47(c)); a do lasso,
% emendada por D49(a), e a Proposition S8.1.
% E1.14 (integrada pela E5e, 2026-10-04): os itens (iii) e (iv) do Corolario 14,
% s' <= 1/4 com s > 1/2, s' <= s/2 com s < 1/2 e a fronteira s = 1/2, estao na
% Proposition S6.1 do supp (prop:slow-ceiling); aqui, uma frase depois da
% proposicao, e a frase de D49(b) revista (D52).

% derivations/08-blocos.tex, Corolario 14 (numeracao global); D47(c), D49(b).

\begin{proposition}[Slow rate, no design condition]
\label{prop:slow}
Let Assumptions~\ref{ass:sample}, \ref{ass:basis} and \ref{ass:besov} hold,
together with the upper bounds in Assumptions~\ref{ass:X} and \ref{ass:U}
($\max_{\cv}|X_{\cv}| \le B_X$, $\lmaxx\{\E(\bX\bX\T\mid\bU)\} \le \kappa_2$ and
$p_{\bU} \le C_U$), with $\Amat$ of rank $p$, the weights as in
Theorem~\ref{thm:main}, $\pen = \lamG$ as in \eqref{eq:lambda} with any constant
$\Lambda > \kappa_2C_U$, $2^{J_n} \le 2n$ and
$pq2^{J_n}\log(pq2^{J_n})/n \to 0$. Then
\[
  \normn{\widehat{f}-f}^{2} = O_p\bigl\{\lamG\nGw{\thetastar} +
  \mathcal{B}_{J_n}^{2} + p/n\bigr\},
\]
with $\lamG\nGw{\thetastar} \le \bar\varrho\,\lamGone\nGone{\thetastar}
\le \bar\varrho\,\lamGone\norm{\thetastar}_{1}$, where $\nGone{\thetastar} = \sum_{G}\norm{\thetastar_{G}}_{2}$ and
$(\lamGone)^{2}$ is of order $\sigma^{2}\log n/n$. Write
$\vartheta_\pi = \min(1-1/\pi,\,1/2)$.
\begin{enumerate}
\item[(i)] If $s > 1/2$ and $\seff > 1/4$, then $\nGone{\thetastar} = O(1)$
and $\normn{\widehat{f}-f}^{2} = O_p\{(\log n/n)^{1/2}\}$ for every
$J_n \ge \log_2 n/(4\seff)$ that satisfies the conditions above, and one
does.
\item[(ii)] If $s < 1/2$ and $\seff > s/2$, then
$\nGone{\thetastar} = O\{1 + \bn^{-\vartheta_\pi}2^{J_n(1/2-s)}\}$, and the
choice $2^{J_n} \asymp \{n/\bn^{(2/\pi-1)_{+}}\}^{1/(1-2s+4\seff)}$ gives
\[
  \normn{\widehat{f}-f}^{2}
  = O_p\bigl\{n^{-2\seff/(1-2s+4\seff)}
  (\log n)^{(2/\pi-1)_{+}2\seff/(1-2s+4\seff)}\bigr\} ,
\]
which, when $\pi \ge 2$, is $n^{-2s/(2s+1)}$, the rate of
Theorems~\ref{thm:main} and \ref{thm:sieve}, with no logarithmic factor.
\end{enumerate}
The proof uses no lower bound on $\Gramhat$, and no property of $\Gramhat$
other than  an upper bound on
$\lmaxx(\Gramhat)$.
\end{proposition}


The bound on $\nGone{\thetastar}$ is Lemma~S6.2 in the
Supplementary Material. Through the chain to $\norm{\thetastar}_{1}$,
Proposition~\ref{prop:slow} reproduces the slow rate of the lasso
(Proposition~S8.1), with the constant $\bar\varrho$: a chunk of up to $2\bn-1$
coefficients enters $\lamGone$ at the order of the union bound over the chunks.
What the chain throws away is the count: at the fine levels
$\nGone{\thetastar}$ is smaller than $\norm{\thetastar}_{1}$ by the factor
$\bn^{-\vartheta_\pi}$, which in (ii) cancels the factor $\bn^{1/2}$ that the
size of the chunks puts in $\lamGone$, entirely when $\pi \ge 2$ and in part
when $\pi < 2$; in (i) the coarse chunk dominates both norms and nothing
changes. No lower bound on $\Gramhat$ enters: the only property of $\Gramhat$
that the proof uses is  an upper bound on
$\lmaxx(\Gramhat)$, here $\lmaxx(\Gramhat) \le \Lambda$, which makes the
deterministic level \eqref{eq:lambda} admissible and holds with probability
tending to one under the upper bounds of Assumptions~\ref{ass:X}
and~\ref{ass:U} alone, so that neither $c_U$ nor $\kappa_1$ is needed. The
conditions $\seff > 1/4$ in (i) and $\seff > s/2$ in (ii) bite only when
$\pi < 2$, and make the resolution of the balance admissible.
% E5e (D52): os pares que (i) e (ii) excluem e a fronteira s = 1/2 vao ao supp.
The pairs that these conditions leave out, $\seff \le 1/4$ with $s > 1/2$ and
$\seff \le s/2$ with $s < 1/2$, all with $\pi < 2$, and the boundary $s = 1/2$
call for a resolution above $n/\log n$, where the column norms and
$\lmaxx(\Gramhat)$ grow with $2^{J_n}/n$; Proposition~S6.1 in the
Supplementary Material covers them with the level inflated to
$\lamGp = (\mu^{\Gcal}_{J_n})^{1/2}\lamG$, where
$\mu^{\Gcal}_{J} = 1 + c_{\mu}2^{J}\log(2D)/n$ for a constant $c_{\mu}$ given
there, so that the first display of Proposition~\ref{prop:slow}, with $\lamGp$
in place of $\lamG$, holds for every $\seff > 0$. There, too, the only
property of $\Gramhat$ that the proof uses is an upper bound,
$\lmaxx(\Gramhat) \le \Lambda\mu^{\Gcal}_{J_n}$, which holds with probability
tending to one with no condition on $J_n$ and, again, without $c_U$ or
$\kappa_1$.
% E5f (E1.15; D54): o Corolario 15 de derivations/08-blocos.tex, que e a
% Proposition S6.2 do supp; a versao do lasso (Proposicao 8 do 05) e a
% Proposition S8.2.
When $s < 1/2$ the slow rate improves further. Read at the comparator that
keeps the chunks with $\norm{\thetastar_{G}}_{2} > \pen w_{G}$ and sets the
others to zero, Theorem~\ref{thm:oracle}(i) charges the penalty only for the
chunks kept and leaves the energy of the others to the bias, where an upper
bound on $\lmaxx(\Gram)$ controls it; with $\seff > s/(2s+1)$ and $J_n$ in the
window \eqref{eq:window}, this gives the prediction rate of
Theorem~\ref{thm:main} with no design condition (Proposition~S6.2 in the
Supplementary Material). When $s \ge 1/2$ no comparator does better than the
order $\lamGone \asymp (\log n/n)^{1/2}$ of part (i): a coarse coefficient of
fixed size costs its product with $\lamGone$ in the penalty or its square in
the bias. Neither reading reaches the components, which need the design
condition (Corollary~\ref{cor:components}).

The second is the approximation lemma that produces Theorem~\ref{thm:main}.

For $\btheta \in \R^{d}$ and $\eta > 0$, let
$\Rid(\btheta;\eta) = \sum_{G\in\Gcal}\min(\norm{\btheta_{G}}_{2}^{2},\eta)$
be the ideal risk over chunks: the risk of an oracle that keeps a chunk when
its energy exceeds $\eta$, at the price $\eta$, and drops it otherwise.

% derivations/05-taxas.tex, Lema 9 (numeracao global). Substituido em k = 4 pelo
% Lema 15 de derivations/08-blocos.tex; o do lasso e o Lemma S8.1.

% derivations/08-blocos.tex, Lema 15 (numeracao global).

\begin{lemma}[Besov regularity bounds the ideal risk over chunks]
\label{lem:weak}
Under Assumptions~\ref{ass:basis} and \ref{ass:besov}, let $b \ge 1$ and let
$\Gcal$ be a partition in which, in each block, the levels with $2^{\lev} < b$
form at most one chunk and every level with $2^{\lev} \ge b$ is cut into
chunks contained in it with at least $b$ coefficients each, as the partition of
Section~\ref{sec:estimator} is. Then, for every $J$ and every $\eta > 0$,
\[
  \Rid(\thetastar;\eta) \;\le\; pq\Bigl\{\AG\,C_g^{\tau}
  \Bigl(\frac{\eta}{b}\Bigr)^{2s/(2s+1)}b^{(2/\pi-1)_{+}/(2s+1)} + \eta\Bigr\},
  \qquad \tau = \frac{1}{s+1/2},
\]
with $\AG = 2 + \{1-2^{1-\pi(s+1/2)}\}^{-1}$ if $\pi \le 2$ and
$\AG = 2 + (1-2^{-2s})^{-1}$ if $\pi \ge 2$.
\end{lemma}



The exponent is the one of Lemma~4 in the supplementary material
of \citet{Klopp-Pensky-2015}, but the assumption is weaker: their class is a
weighted $\ell_\nu$ ball over all the coefficients, which implies
Assumption~\ref{ass:besov} level by level, and Assumption~\ref{ass:besov} does
not need the sum over the levels. The coarse chunk costs at most one $\eta$
per block.




Theorem~\ref{thm:main} follows by reading
Theorem~\ref{thm:oracle}(iii) not at the oracle but at the comparator that
keeps the chunks of $\thetastar$ whose squared norm exceeds a level $\eta$ and
sets the others to zero, which the theorem allows because it holds for every
comparator at once. The estimation term then counts $\eta$ for each chunk
kept, and the truncation leaves the energy of the others, so that the bound is
of the order of $\Rid(\thetastar;\eta)$ (Corollary~S5.1 in the Supplementary
Material), with $\eta$ of the order of $\max_{G}(\pen w_{G})^{2}$, that is of
$\bn/n$. Lemma~\ref{lem:weak} with $b = \bn$ turns this into
$n^{-2s/(2s+1)}\bn^{(2/\pi-1)_{+}/(2s+1)}$: the chunk pays $\bn$ in the level
$\eta$ and returns it through $(\eta/b)^{2s/(2s+1)}$, so that the cost is of
order $1/n$ per coefficient, against $\log n/n$ for the lasso, and the
logarithm of the union bound over the columns is absorbed by the size of the
chunks, which is why they must have at least of the order of $\log n$
coefficients. When $\pi < 2$ the bound charges
$\bn^{(2/\pi-1)/(2s+1)}$ to sequences whose large coefficients sit one per
chunk; a signal spread over a level pays nothing for the chunks. The two
inequalities that define the window \eqref{eq:window} are that the
approximation bias not dominate and that the requirement of
Proposition~\ref{prop:gram}(ii) hold. Nothing in this argument says that the
estimator selects the chunks kept by the comparator: no minimal signal or
irrepresentability condition is used.
The recovery that matters for the model is not of coordinates
but of blocks, and it is taken up next.

% derivations/06-selecao-limiar.tex, Lema 13, Hipotese S e Corolario 8
% (numeracao global); D32. O Lema 13 e a prova estao no supp (Secao S7).

\subsection{Which modulators act on which coefficients}
\label{sec:selection}

The structure of the model is the set of blocks whose component does not
vanish,
\[
  \mathcal{S} \;=\; \bigl\{(\cv,\md) :\; \nu_{\cv\md} > 0\bigr\},
  \qquad
  \nu_{\cv\md} = \bigl\lVert \gcomp{\cv}{\md} \bigr\rVert_{L_2[0,1]},
\]
and to know it is to know which modulators act on which coefficients. The
 estimator \eqref{eq:estimator} is not claimed to recover it: its
penalty removes chunks and not blocks, under Assumption~\ref{ass:besov} the
oracle is dense, and nothing prevents the fit from keeping nonzero
coefficients in blocks whose component is zero. What the preceding results do give is that every block is
estimated at the rate of Corollary~\ref{cor:components}, and that is enough to
recover $\mathcal{S}$ in a second step, by a threshold on the norms of the fitted
components. The procedure is estimation followed by a threshold. By the
orthonormality of the basis, the norm of a fitted component is the Euclidean
norm of its block of coefficients,
\[
  \widehat{\nu}_{\cv\md}
  = \bigl\lVert \widehat{g}_{\cv\md} \bigr\rVert_{L_2[0,1]}
  = \norm{\thetahat_{\cv\md,\cdot}}_{2},
  \qquad
  \widehat{\mathcal{S}}(t) = \bigl\{(\cv,\md) :\; \widehat{\nu}_{\cv\md} > t\bigr\},
\]
so that the statistic costs nothing once the estimator has been fit. Write
 $\rhoG = n^{-\seff/(2\seff+1)}$ for the square root of the rate of
Theorem~\ref{thm:sieve}.

\begin{assumption}[Block separation]
\label{ass:sep}
The components may depend on $n$, within the class of
Assumption~\ref{ass:besov} and with its constants fixed, and so may $\mathcal{S}$; the
smallest active block, $\nu_{\min,n} = \min_{(\cv,\md)\in \mathcal{S}} \nu_{\cv\md}$, with
$\nu_{\min,n} = \infty$ when $\mathcal{S}$ is empty, satisfies  $\nu_{\min,n}/\rhoG \to \infty$.
\end{assumption}

\begin{corollary}[Block structure by thresholding]
\label{cor:threshold}
Let the assumptions of Theorem~\ref{thm:sieve} hold, with $J_n$ and  $\pen = \lamG$ as
there.
\begin{enumerate}
\item[(i)] (Screening.) If the thresholds satisfy  $t_n/\rhoG \to \infty$,
then $\Prob\{\widehat{\mathcal{S}}(t_n) \subseteq \mathcal{S}\} \to 1$.
\item[(ii)] (Recovery.) If, in addition, Assumption~\ref{ass:sep} holds and
$\limsup_n t_n/\nu_{\min,n} < 1$, then $\Prob\{\widehat{\mathcal{S}}(t_n) = \mathcal{S}\} \to 1$.
Such thresholds always exist; one is  $t_n = (\rhoG\nu_{\min,n})^{1/2}$.
\item[(iii)] (A window at fixed $n$.) With $\alpha$ the level in
\eqref{eq:lambda}, for every $n$ and every $\alpha' \in (0,1)$ there is an
event, of probability at least
$1 - \alpha - \alpha' - \omega_n$ with $\omega_n \to 0$ not depending on the
components, on which $\widehat{\mathcal{S}}(t) = \mathcal{S}$ for every $t$
with  $\bar{\Delta}^{\Gcal}_n \le t < \nu_{\min,n} - \bar{\Delta}^{\Gcal}_n$; here
 $\bar{\Delta}^{\Gcal}_n \le c(\alpha,\alpha')\rhoG$ is an explicit bound on
$\max_{(\cv,\md)}\lVert \widehat{g}_{\cv\md} - \gcomp{\cv}{\md}
\rVert_{L_2[0,1]}$, given in the proof (Section~S7 of the Supplementary
Material).
\end{enumerate}
\end{corollary}

Parts (i) and (ii) hold also at the resolution of
Theorem~\ref{thm:main}, with  $\rhoG$ replaced by
 $\rhoGt = n^{-s/(2s+1)}(\log n)^{(2/\pi-1)_{+}/\{2(2s+1)\}}$, the square root of the rate of that
theorem, both in the condition on $t_n$ and in Assumption~\ref{ass:sep}.
Since  $\rhoGt \le \rhoG$ for large $n$, with a gap in the exponent when $\pi < 2$,
the separation required is weaker exactly where the components are spatially
inhomogeneous. Part (iii) is stated at the resolution of
Theorem~\ref{thm:sieve}, where its bound is explicit. The componentwise rate of
the lasso at its own resolution is $(\log n/n)^{\seff/(2\seff+1)}$
(Corollary~S8.1 in the Supplementary Material), so that the separation
required here is weaker than for the lasso by the factor
$(\log n)^{\seff/(2\seff+1)}$.

Part (i) uses no separation: removing the null blocks is free, and
Assumption~\ref{ass:sep} is spent entirely on not losing a small active block.
Under a fixed target $\nu_{\min,n}$ is a positive constant and the assumption holds
automatically, because  $\rhoG \to 0$; it carries content in a triangular
array, where it says how fast the smallest active block may shrink. It bears
on the components and not on $(\bX,\bU)$, so that no design condition beyond
Proposition~\ref{prop:gram} and no irrepresentability condition enters; a
separation condition on the components plays the same role in the selection
of variables in nonparametric additive models \citep{Huang-Horowitz-Wei-2010}.
The proof rests on a deterministic sandwich (Lemma~S7.1 in the Supplementary
Material): if every component is estimated within $\bar{\Delta}$ in $L_2[0,1]$, then
$\{\nu_{\cv\md} > t + \bar{\Delta}\} \subseteq \widehat{\mathcal{S}}(t) \subseteq
\{\nu_{\cv\md} > t - \bar{\Delta}\}$, so that every $t \in [\bar{\Delta}, \nu_{\min,n} - \bar{\Delta})$ recovers $\mathcal{S}$,
and that window cannot be widened by the argument. Two things are left out.
The window involves $\sigma$, the design constant and $\seff$, so the
corollary does not say how to choose $t_n$ from the data (Section~\ref{sec:threshold}
uses the one standard error rule of cross-validation, for which no theorem is
given here); and it says nothing
about which wavelets are active inside a block.

% derivations/06-selecao-limiar.tex, Lema 16 e Corolario 13 (numeracao global;
% adendo de E1.12); D45. O Lema 16 e a prova estao no supp (Secao S7).
Thresholding is also the estimator used for prediction
(Section~\ref{sec:threshold}). Write $\thetahat^{(t)}$ for $\thetahat$ with the
blocks outside $\widehat{\mathcal{S}}(t)$ set to zero, and
$\widehat{g}^{(t)}_{\cv\md}$ and $\widehat{f}^{(t)} = \Amat\chat +
\Bmat\thetahat^{(t)}$ for the resulting components and fit, so that the levels
and the blocks kept are those of the fit. Setting to zero an active block
whose fitted norm is at most $t$ costs at most its norm, which is at most
$t + \lVert\widehat{g}_{\cv\md} - \gcomp{\cv}{\md}\rVert_{L_2[0,1]}$, setting to
zero a null block only reduces the error, and the blocks kept do not change
(Lemma~S7.2 in the Supplementary Material). With the rates of the fit this
gives the risk of the thresholded estimator.

\begin{corollary}[Risk of the thresholded estimator]
\label{cor:thrisk}
Let $(\chat,\thetahat)$ be either the fit of Theorem~\ref{thm:sieve}, with
$r_n = \rhoG$, or the lasso at the resolution of Theorem~S8.2 in the
Supplementary Material, with $r_n = (\log n/n)^{\seff/(2\seff+1)}$, and let
$t_n \ge 0$.
\begin{enumerate}
\item[(i)] With no further assumption,
$\sum_{\cv,\md}\lVert\widehat{g}^{(t_n)}_{\cv\md}-\gcomp{\cv}{\md}\rVert_{L_2[0,1]}^{2}
= O_p(r_n^{2} + t_n^{2})$ and
$\normn{\widehat{f}^{(t_n)}-f}^{2} = O_p(r_n^{2} + t_n^{2})$.
\item[(ii)] If Assumption~\ref{ass:sep} holds with $r_n$ in place of $\rhoG$
and $\limsup_n t_n/\nu_{\min,n} < 1$, then with probability tending to one
$\mathcal{S} \subseteq \widehat{\mathcal{S}}(t_n)$,
\[
  \sum_{\cv,\md}\lVert\widehat{g}^{(t_n)}_{\cv\md}-\gcomp{\cv}{\md}\rVert_{L_2[0,1]}^{2}
  \;\le\; \sum_{\cv,\md}\lVert\widehat{g}_{\cv\md}-\gcomp{\cv}{\md}\rVert_{L_2[0,1]}^{2}
\]
and
\[
  \normn{\widehat{f}^{(t_n)}-f}^{2} \;\le\;
  \frac{4\Lambda}{\gamma}\normn{\widehat{f}-f_{J_n}}^{2} + 2\normn{\bb}^{2},
\]
so that the thresholded estimator keeps the rate $r_n^{2}$ of the fit. If
moreover $t_n/r_n \to \infty$, then
$\widehat{\mathcal{S}}(t_n) = \mathcal{S}$ with probability tending to one,
and the null blocks no longer contribute to the error of the components.
\end{enumerate}
\end{corollary}

Under a fixed target Assumption~\ref{ass:sep} holds automatically, and (ii)
holds for every $t_n \to 0$. The factor $4\Lambda/\gamma$ in (ii) bounds the
condition number of $\Gramhat$, through which the bound passes from
prediction to coefficients and back; it is a price of the argument, and the
corollary says what any threshold costs, not which threshold a data driven
rule selects.


%%%%%%%%%%%%%%%%%%%%%%%%%%%%%%%%%%%%%%%%%%%%%%%%%%%%%%%%%%%%%%%%%%%%%%%%%%%%%%
